\chapter{Introduction}

\section{Motivation}
Modern network infrastructures generate vast amounts of traffic, complicating the task of distinguishing legitimate communication from malicious activity. Network Intrusion Detection Systems (NIDS) rely on statistical patterns, flow characteristics, and behavioral modeling to flag potential threats. As machine learning techniques mature, they offer promising capabilities for managing this complexity. However, standard supervised learning formulates intrusion detection as a symmetric classification problem, often minimizing overall error rather than operational risk. In cybersecurity, the cost of missing a severe attack (a false negative) frequently outweighs the cost of investigating a false alert (a false positive). Reinforcement Learning (RL) provides a natural framework for asymmetric, cost-sensitive decision-making by allowing defenders to optimize for specific reward structures. Exploring how RL algorithms navigate these trade-offs on flow-level representations can provide valuable insights into risk-aware automated defense.

\section{Problem Statement}
Applying RL to intrusion detection involves significant methodological challenges. First, high-dimensional network flow data suffers from distribution shifts, class imbalance, and artifacts intrinsic to public datasets. Second, RL models often require extensive data and careful tuning, making it difficult to assess whether they offer empirical advantages over robust, well-understood supervised baselines like Random Forests. Furthermore, many existing studies evaluate models on isolated, randomized partitions of legacy datasets, failing to expose data leakage or to measure generalization under strict temporal splits. The problem is therefore not merely applying an RL algorithm to a network dataset, but rather structuring the problem so that the model's cost-sensitive decisions can be evaluated rigorously, safely, and transparently against established benchmarks.

\section{Proposed Approach}
This thesis formulates the defensive task as an offline, dataset-as-environment reinforcement learning problem. Network traffic is processed into a fixed, 76-feature canonical flow schema. To manage missing or heterogeneous data, a 76-value missingness mask is appended, producing a 152-dimensional observation vector for the agent. The decision task is mapped to a binary action space: \texttt{PERMIT} (classify as benign) or \texttt{BLOCK} (classify as attack). A Quantile Regression Deep Q-Network (QRDQN) agent is trained within a Gymnasium-compatible environment that explicitly penalizes false negatives more severely than false positives. The CICIDS2017 dataset serves as the primary internal benchmark, strictly processed to exclude leakage-prone fields such as IP addresses and absolute timestamps. Performance is then benchmarked against a Random Forest baseline under progressively stricter validation constraints.

\section{Contribution of the Thesis}
This work contributes a reproducible, leakage-aware, and cost-sensitive experimental prototype for evaluating RL-based binary security decisions on flow records. It explicitly defines a canonical feature contract that separates dataset artifacts from the model's observation space. Furthermore, the thesis implements a multi-stage validation ladder---ranging from simple random splits to leave-one-CSV-out evaluations---providing a structured methodology for identifying domain shifts and assessing model robustness. Finally, it presents a rigorous empirical comparison between QRDQN and supervised baselines, offering evidence on the data-efficiency and trade-offs of the RL approach without overclaiming deployment readiness.

\section{Scope and Limitations}
The scope of this research is strictly confined to offline, tabular flow-level evaluation. The \texttt{PERMIT} and \texttt{BLOCK} actions are experimental binary classifications; the system does not perform live firewall blocking, packet dropping, or inline intrusion prevention. Furthermore, the environment does not simulate an active adversary whose subsequent actions change in response to the defender's choices. Consequently, this system is an analytical prototype, not a production-ready NIDS or an autonomous cyber-defense control plane. CICIDS2017 is utilized as an internal benchmark for controlled experimentation, not as definitive proof of real-world generalization, with external laboratory validation prioritized as a subsequent, stricter evaluation tier.

\section{Structure of the Document}
The remainder of this document is organized as follows. Chapter 2 details the specific objectives, scope, and non-goals of the thesis. Chapter 3 reviews the State of the Art, surveying flow-based NIDS, public datasets, supervised baselines, and RL methodologies, while discussing common evaluation pitfalls. Subsequent chapters outline the methodological design, the environment architecture, the experimental setup, and the final empirical results.
